\begin{scanpage}{0211}{193}
\noindent in it gives us the field radiated by the semi-infinite conductor $(z_1,\infty)$. The wave radiated by the conductor's end is the desired secondary edge wave and may be represented in the form
\[E_\vartheta^{(2)}(z_1)=H_\varphi^{(2)}(z_1)=\frac{\mathcal E^{(2)}(z_1)}{2\sin\vartheta\ln\dfrac{2i}{\gamma ka\sin\vartheta}}\frac{e^{ikR}}R e^{-ikz_1\cos\vartheta},\tag{30.34}\]
where
\[
\begin{aligned}
\mathcal E^{(2)}(z_1)={}&-\frac{4iE e^{iw_0z_2+ikL}}{k\sin\vartheta_0(\cos\vartheta+\cos\vartheta_0)\ln\dfrac{2i}{\gamma ka\sin\vartheta_0}}\times\\
&\times\left[\sin^2\frac{\vartheta_0}2\sin^2\frac\vartheta2\ln\frac{i}{\gamma ka\sin\dfrac{\vartheta_0}2}\psi_-(L)\right.\\
&\qquad\left.-\cos^2\frac{\vartheta_0}2\cos^2\frac\vartheta2\ln\frac{i}{\gamma ka\cos\dfrac\vartheta2}\psi_+^0(L)\right].
\end{aligned}\tag{30.35}
\]
Otherwise, this expression may be written directly by replacing, in Equations (30.30) and (30.31), $z_2$ by $z_1$, $\vartheta$ by $\pi-\vartheta$ and $\vartheta_0$ by $\pi-\vartheta_0$.
\sectiontitle{31}{Multiple Diffraction of Edge Waves}
The secondary waves (30.30) and (30.34) which we found are the waves diverging from the ends of the semi-infinite conductors $(-\infty,z_2)$ and $(z_1,\infty)$. If one excites an infinite single-wire line by the external field
\[E_z^e=\mathcal E_2^{(2)}\delta(z-z_2),\tag{31.01}\]
where
\[\mathcal E_2^{(2)}=\left.\mathcal E^{(2)}(z_2)\right|_{\vartheta=\pi}
=E\frac{2i\ln\dfrac{i}{\gamma ka}\psi_-^{(0)}(L)}{k\sin\vartheta_0\ln\dfrac{2i}{\gamma ka\sin\vartheta_0}}e^{ikL-ikz_1\cos\vartheta_0},\tag{31.02}\]
then a spherical wave arises which with $\vartheta\approx\pi$ coincides with wave (30.30). With the excitation of an infinite line by the concentrated emf
\end{scanpage}
